\documentclass{article}
\usepackage{amsmath, amssymb, amsthm}
\usepackage{hyperref}
\usepackage{geometry}
\geometry{margin=1in}

\title{Erd\H{o}s Problem \#886}
\date{Last updated: 2025-08-31}
\author{Source: erdosproblems.com}

\begin{document}
\maketitle

\noindent\textbf{Status:} OPEN \quad \textbf{No prize}

\noindent\textbf{Tags:} number theory, divisors

\medskip

\noindent\textbf{Problem Statement:}

Let $\epsilon&#62;0$. Is it true that, for all large $n$, the number of divisors of $n$ in $(n^{1/2},n^{1/2}+n^{1/2-\epsilon})$ is $O_\epsilon(1)$?




Erd\H{o}s attributes this conjecture to Ruzsa. Erd\H{o}s and Rosenfeld \cite{ErRo97} proved that there are infinitely many $n$ such that there are four divisors of $n$ in $(n^{1/2},n^{1/2}+16n^{1/4})$. They also proved that, for any constant $C&gt;0$, all large $n$ have at most $1+C^2$ many divisors in\[[n^{1/2}, n^{1/2}+Cn^{1/4}].\]See also [887] .




References

[ErRo97] Erd\H{o}s, Paul and Rosenfeld, Moshe, The factor-difference set of integers . Acta Arith. (1997), 353--359.


\bigskip
\noindent\small{Source: \url{https://www.erdosproblems.com/886}}
\end{document}
